% !TEX program = pdflatex
\documentclass[11pt,a4paper]{article}

% =========================================================================
% VIABILIDAD Y BUSINESS CASE
% Proyecto: CasaScore (AlquilaPuno)
% Gestión de Proyectos TICS (EST315) -- Semestre 2026-II -- Team Sobrevivientes
% =========================================================================

% --- PAQUETES DE CONFIGURACIÓN ---
\usepackage[utf8]{inputenc}
\usepackage[T1]{fontenc}
\usepackage[spanish,es-tabla]{babel}
% Compatibilidad con babel-spanish reciente (flechas de TikZ y signo de porcentaje)
\AtBeginDocument{\shorthandoff{<>}\spanishplainpercent}
\usepackage{geometry}
\geometry{
    top=3.6cm,
    bottom=2.3cm,
    left=2.2cm,
    right=2.2cm,
    % PARTE MODIFICABLE: headheight debe ser mayor que la altura del logo más alto (2.56cm = 73pt).
    headheight=80pt,
    headsep=15pt,
    footskip=20pt
}

\usepackage{graphicx}
\usepackage{fancyhdr}
\usepackage[table]{xcolor}
\usepackage{titlesec}
\usepackage{booktabs}
\usepackage{tabularx}
\usepackage{array}
\usepackage{amsmath}
\usepackage{enumitem}
\usepackage{microtype}
\usepackage{amssymb}
\usepackage{hyperref}

% Logos: funciona a 2 o 3 niveles de profundidad dentro de docs/latex/
\graphicspath{{../../assets/img/}{../../../assets/img/}}

% --- PALETA INSTITUCIONAL FIEI / UNA PUNO ---
\definecolor{AzulUNA}{RGB}{10, 45, 90}
\definecolor{DoradoFIEI}{RGB}{180, 130, 20}
\definecolor{GrisOscuro}{RGB}{50, 50, 50}
\definecolor{AzulClaro}{RGB}{228, 236, 246}

% --- CONFIGURACIÓN DE LISTAS ---
\setlist{topsep=2pt, itemsep=2pt, parsep=1pt}

% --- COLUMNAS DE TABLA ---
\newcolumntype{L}{>{\raggedright\arraybackslash}X}
\newcolumntype{P}[1]{>{\raggedright\arraybackslash}p{#1}}
\newcolumntype{C}[1]{>{\centering\arraybackslash}p{#1}}
\newcolumntype{Q}[1]{>{\raggedleft\arraybackslash}p{#1}}
\newcolumntype{R}{>{\raggedleft\arraybackslash}X}
\renewcommand{\arraystretch}{1.25}

% --- FORMATO DE ENLACES ---
\hypersetup{
    colorlinks=true,
    linkcolor=AzulUNA,
    urlcolor=AzulUNA,
    citecolor=AzulUNA,
    pdftitle={Viabilidad y Business Case -- CasaScore (AlquilaPuno)},
    pdfauthor={Team Sobrevivientes}
}

% --- FORMATO DE TÍTULOS ---
\titleformat{\section}{\color{AzulUNA}\normalfont\large\bfseries}{\thesection.}{0.7em}{}[\color{AzulUNA}\titlerule]
\titleformat{\subsection}{\color{AzulUNA}\normalfont\normalsize\bfseries}{\thesubsection.}{0.5em}{}
\titlespacing*{\section}{0pt}{9pt}{4pt}
\titlespacing*{\subsection}{0pt}{6pt}{2pt}

% --- ENCABEZADO Y PIE DE PÁGINA (fancyhdr) ---
\pagestyle{fancy}
\fancyhf{}

% ===== TAMAÑO DE LOS LOGOS (PARTE MODIFICABLE) =====
% logo_fiei.png  : height=1.45cm (lado izquierdo)
% logo_equipo.png: height=2.56cm (lado derecho; 25 % más grande porque la imagen tiene más borde blanco)
% Para cambiar el tamaño, modifique solo el valor height=... de cada \includegraphics.
% Si el logo más alto supera 2.6cm, aumente también headheight y top en la configuración de página.
\fancyhead[L]{%
    \includegraphics[height=1.45cm,keepaspectratio]{logo_fiei.png}%
}
\fancyhead[C]{%
    \scriptsize \textbf{\textcolor{AzulUNA}{UNIVERSIDAD NACIONAL DEL ALTIPLANO DE PUNO}}\\[2pt]
    \tiny \textcolor{GrisOscuro}{FACULTAD DE INGENIERÍA ESTADÍSTICA E INFORMÁTICA}\\[2pt]
    \scriptsize \textbf{\textcolor{AzulUNA}{Gestión de Proyectos TICS (EST315) \textbar\ Team Sobrevivientes}}\\[2pt]
    \footnotesize \textbf{\textcolor{DoradoFIEI}{VIABILIDAD Y BUSINESS CASE}}%
}
\fancyhead[R]{%
    \includegraphics[height=2.56cm,keepaspectratio]{logo_equipo.png}%
}

\renewcommand{\headrulewidth}{0pt}
\renewcommand{\headrule}{}

\fancyfoot[L]{\scriptsize \textcolor{GrisOscuro}{CasaScore (AlquilaPuno)}}
\fancyfoot[C]{\small \thepage}
\fancyfoot[R]{\scriptsize \textcolor{GrisOscuro}{Puno, Perú -- Semestre 2026-II}}
\renewcommand{\footrulewidth}{0.4pt}

\begin{document}

\begin{center}
    {\Large \textbf{\textcolor{AzulUNA}{VIABILIDAD Y BUSINESS CASE}}}\\[3pt]
    {\normalsize \textbf{CasaScore (AlquilaPuno)}}\\[4pt]
    \footnotesize \textbf{Fecha:} jueves 10/09/2026 (semana 3) \quad \textbar\quad \textbf{Versión:} 1.0 \quad \textbar\quad \textbf{Fase:} Inicio
\end{center}

\vspace{0.1cm}

\section{Datos Generales}
\noindent\begin{tabularx}{\textwidth}{@{}P{4.4cm} L@{}}
\toprule
\textbf{Proyecto:} & CasaScore (aplicación AlquilaPuno) \\
\textbf{Responsable:} & Ruth Karina Apaza Solis -- Analista de Planificación \\
\textbf{Sponsor:} & Team Sobrevivientes, dueño y patrocinador del proyecto \\
\textbf{Tipo de estimación (PM4DEV):} & Aproximada, con un rango de $-25\,\%$ a $+75\,\%$ \\
\textbf{Horizonte de evaluación:} & Semestre 2026-II (inversión) y años 1 a 3 (operación); tasa de descuento de 10\,\% \\
\textbf{Moneda:} & Soles (S/) \\
\bottomrule
\end{tabularx}

\section{Matriz de Viabilidad}
{\small
\noindent\begin{tabularx}{\textwidth}{@{}P{2.3cm} L P{2.5cm}@{}}
\toprule
\rowcolor{AzulClaro}
\textbf{Dimensión} & \textbf{Análisis} & \textbf{Dictamen} \\
\midrule
\textbf{Técnica} & El equipo domina las herramientas (Expo, React Native, TypeScript y SQL). El motor de puntaje usa normalización y suma ponderada, técnicas de estadística aplicada conocidas por el equipo. Mapas abiertos (OpenStreetMap) y base de datos en la nube (Supabase) con capa gratuita. & Viable \\
\addlinespace[2pt]
\textbf{Económica} & Costo valorizado del semestre: S/ 11,108.07, de los cuales S/ 2,507.44 son gasto en efectivo, cubierto en partes iguales por los 4 integrantes. Los ingresos de la app no cubren la operación en los 3 primeros años, pero el beneficio para los estudiantes es alto (sección 5). & Viable con condiciones: buscar fondos en concursos de la UNA Puno y ferias \\
\addlinespace[2pt]
\textbf{Operativa} & El equipo opera la app después del semestre como dueño del producto. Los propietarios y estudiantes son accesibles en el entorno de la UNA Puno para el piloto. & Viable \\
\addlinespace[2pt]
\textbf{Legal} & Datos personales tratados según la Ley N.\textsuperscript{o} 29733, con consentimiento y política de privacidad. Los cobros desde el año 1 se emiten con el RUC de Clyde Neil Paricahua Pari como persona natural con negocio. La app no intermedia el pago del alquiler ni firma contratos. & Viable \\
\addlinespace[2pt]
\textbf{Temporal} & MVP de la Unidad I en 5 sprints (semanas 4 a 8) y MVP final en la semana 16, con 8 h semanales por integrante. & Viable \\
\bottomrule
\end{tabularx}
}

\section{Costos del Proyecto (semestre 2026-II)}

\subsection{Parámetros}
Tarifa por hora $=$ sueldo de mercado $/$ 192 h $\times$ 1.09 (EsSalud). Dedicación: 8 h por semana por integrante durante 16 semanas (128 h). Se prevén 6 jornadas nocturnas antes de las entregas, con 4 horas extra por persona (sobretasa de 25\,\% y 35\,\%). Equipos depreciados al 25\,\% anual. Tipo de cambio: S/ 3.436 por dólar.

\subsection{Estimación por rubros}
{\small
\noindent\begin{tabularx}{\textwidth}{@{}P{3.6cm} L Q{2.0cm}@{}}
\toprule
\rowcolor{AzulClaro}
\textbf{Rubro} & \textbf{Detalle} & \textbf{S/} \\
\midrule
Personal regular & 512 h del equipo: PM S/ 14.19/h, Analista S/ 11.35/h, Técnico S/ 12.49/h y Gestora S/ 10.22/h & 6,176.00 \\
Horas extra & 6 jornadas nocturnas $\times$ 4 h $\times$ 4 integrantes & 1,505.40 \\
Equipos & Depreciación de 4 laptops y 1 celular de pruebas & 919.23 \\
Alimentación & Dietas nocturnas, refrigerios y almuerzos de campo & 1,120.00 \\
Productos de salud & Gotas para los ojos, multivitamínico, analgésico y bloqueador solar & 314.00 \\
Servicios & Electricidad, internet del hogar y datos móviles & 540.54 \\
Transporte & Taxi nocturno y pasajes de trabajo de campo & 252.00 \\
Licencias & Cuenta Google Play Console (US\$ 25) & 85.90 \\
Documentación y piloto & Impresiones, afiches e incentivos para 15 participantes & 195.00 \\
\midrule
\textbf{Estimación de costos} & Rango PM4DEV: S/ 8,331.05 a S/ 19,439.12 & \textbf{11,108.07} \\
(+) Contingencia & Valor esperado de los riesgos identificados & 555.40 \\
\textbf{Línea base de costos} & & \textbf{11,663.47} \\
(+) Reserva de gestión 10\,\% & Imprevistos & 1,166.35 \\
\textbf{Presupuesto total} & & \textbf{12,829.82} \\
\bottomrule
\end{tabularx}
}

\subsection{Financiamiento}
{\small
\noindent\begin{tabularx}{\textwidth}{@{}P{4.2cm} L Q{2.0cm}@{}}
\toprule
\rowcolor{AzulClaro}
\textbf{Fuente} & \textbf{Detalle} & \textbf{S/} \\
\midrule
Aporte de los integrantes & Gasto en efectivo de S/ 2,507.44 dividido en partes iguales: S/ 626.86 por integrante & 2,507.44 \\
Trabajo y equipos propios & Horas del equipo y depreciación de equipos, valorizados & 8,600.63 \\
Concursos de la UNA Puno y ferias de innovación & Postulación con el MVP; los premios no se incluyen en el cálculo & -- \\
\bottomrule
\end{tabularx}
}

\section{Ingresos y Costos de Operación (años 1 a 3)}
Usuarios activos al mes: 500, 1,000 y 2,000. Google Play retiene 15\,\% de las compras dentro de la app.

\vspace{0.1cm}
{\small
\noindent\begin{tabularx}{\textwidth}{@{}P{3.0cm} L Q{1.55cm} Q{1.55cm} Q{1.55cm}@{}}
\toprule
\rowcolor{AzulClaro}
\textbf{Concepto} & \textbf{Supuesto} & \textbf{Año 1} & \textbf{Año 2} & \textbf{Año 3} \\
\midrule
Anuncios recompensados & 2 videos por usuario al mes; US\$ 1.8 por cada 1,000 vistas & 74.22 & 148.44 & 296.87 \\
Versión Premium & 2\,\% de usuarios $\times$ S/ 4.90 & 41.65 & 83.30 & 166.60 \\
Anunciantes locales & 5 por ciudad $\times$ S/ 15 al mes & 900.00 & 1,800.00 & 3,600.00 \\
Suscripción por cuarto & S/ 5 por cuarto publicado al mes; 270, 540 y 1,080 cuartos-mes & 1,147.50 & 2,295.00 & 4,590.00 \\
\textbf{Ingresos} & & \textbf{2,163.37} & \textbf{4,326.74} & \textbf{8,653.47} \\
\midrule
Mantenimiento & 4, 6 y 8 h/semana $\times$ S/ 12.49 & 2,597.92 & 3,896.88 & 5,195.84 \\
Moderación y soporte & 2, 3 y 4 h/semana $\times$ S/ 10.22 & 1,062.88 & 1,594.32 & 2,125.76 \\
Base de datos & Supabase Pro (US\$ 25, 25 y 35 al mes) & 1,030.80 & 1,030.80 & 1,443.12 \\
Dominio y difusión & Dominio web y material por ciudad & 341.23 & 641.23 & 1,241.23 \\
\textbf{Costos de operación} & & \textbf{5,032.83} & \textbf{7,163.23} & \textbf{10,005.95} \\
\bottomrule
\end{tabularx}
}

\section{Evaluación del Business Case}

\subsection{Evaluación financiera (solo ingresos de la app)}
{\small
\noindent\begin{tabularx}{\textwidth}{@{}C{1.2cm} R R R@{}}
\toprule
\rowcolor{AzulClaro}
\textbf{Año} & \textbf{Flujo neto} & \textbf{Descontado (10\,\%)} & \textbf{Acumulado} \\
\midrule
0 & $-$11,108.07 & $-$11,108.07 & $-$11,108.07 \\
1 & $-$2,869.46 & $-$2,608.60 & $-$13,977.53 \\
2 & $-$2,836.49 & $-$2,344.21 & $-$16,814.02 \\
3 & $-$1,352.48 & $-$1,016.14 & $-$18,166.50 \\
\midrule
\multicolumn{2}{@{}l}{\textbf{VAN financiero}} & \textbf{$-$17,077.02} & \\
\bottomrule
\end{tabularx}
}

\vspace{0.1cm}
\noindent\textbf{Punto de equilibrio operativo del año 1:} 945 cuartos-mes publicados.

\subsection{Evaluación económica (ingresos más beneficio para los estudiantes)}
Beneficio por estudiante que encuentra cuarto con la app: S/ 41.64 al año (4 h de búsqueda ahorradas $\times$ S/ 6.41, S/ 7.00 de estafas evitadas y S/ 9.00 de movilidad evitada). Estudiantes beneficiados: 350, 700 y 1,400.

\vspace{0.1cm}
{\small
\noindent\begin{tabularx}{\textwidth}{@{}C{1.2cm} R R R R@{}}
\toprule
\rowcolor{AzulClaro}
\textbf{Año} & \textbf{Flujo financiero} & \textbf{Beneficio social} & \textbf{Flujo económico} & \textbf{Descontado} \\
\midrule
0 & $-$11,108.07 & -- & $-$11,108.07 & $-$11,108.07 \\
1 & $-$2,869.46 & 14,574.00 & 11,704.54 & 10,640.49 \\
2 & $-$2,836.49 & 29,148.00 & 26,311.51 & 21,745.05 \\
3 & $-$1,352.48 & 58,296.00 & 56,943.52 & 42,782.51 \\
\bottomrule
\end{tabularx}
}

\vspace{0.15cm}
{\small
\noindent\begin{tabularx}{\textwidth}{@{}P{4.0cm} L Q{3.1cm}@{}}
\toprule
\rowcolor{AzulClaro}
\textbf{Indicador} & \textbf{Cálculo} & \textbf{Resultado} \\
\midrule
VAN económico & $-$11,108.07 + 10,640.49 + 21,745.05 + 42,782.51 & S/ +64,059.98 \\
ROI (3 años) & (94,959.57 $-$ 11,108.07) / 11,108.07 & 754.9\,\% \\
Payback & 11,108.07 / 11,704.54 & 0.95 años ($\approx$ 11 meses) \\
Relación beneficio/costo & 75,168.05 / 11,108.07 & 6.77 \\
TIR & Tasa que hace VAN = 0 & 166.5\,\% \\
\bottomrule
\end{tabularx}
}

\section{Firmas}
\noindent En señal de conformidad con el presente análisis, los integrantes del equipo firman a continuación:

\vspace{1.5cm}
\begin{center}
\begin{tabular}{c c c}
    \parbox{6.5cm}{\centering \rule{6.0cm}{0.6pt}\\[5pt] \textbf{Clyde Neil Paricahua Pari}\\[2pt] {\footnotesize Project Manager}}
    & \hspace{0.8cm} &
    \parbox{6.5cm}{\centering \rule{6.0cm}{0.6pt}\\[5pt] \textbf{Ruth Karina Apaza Solis}\\[2pt] {\footnotesize Analista de Planificación}}
    \\[2.0cm]
    \parbox{6.5cm}{\centering \rule{6.0cm}{0.6pt}\\[5pt] \textbf{Eliseo Tarqui Ajahuana}\\[2pt] {\footnotesize Especialista Técnico / Arquitectura}}
    & \hspace{0.8cm} &
    \parbox{6.5cm}{\centering \rule{6.0cm}{0.6pt}\\[5pt] \textbf{Viviana Mary Quisbert Quispe}\\[2pt] {\footnotesize Gestora de Comunicación, Riesgos y Documentación}}
\end{tabular}
\end{center}
\end{document}
